\section{Notation}
To avoid confusion, let's explicitly define the notation used in the rest of the text.
\begin{itemize}
\item $ q \in \mathbb{R}^{nq} $ - joint configuration vector
\item $ v \in \mathbb{R}^{nv} $ - joint velocity vector
\item $ a \in \mathbb{R}^{nv} $ - joint acceleration vector
\item $ \tau \in \mathbb{R}^{nv} $ - joint torque vector
\item $ u \in \mathbb{R}^{nu} $ - control input (almost always 
	the velocity command vector $ v_{ cmd } \in \mathbb{R}^{nv}  $)
\item $ T_{ a }^{ b } \in SE (3) $ - element of the special Euclidean group,
	representing the frame $ \left\{ b \right\}   $ as seen in frame $ \left\{ a \right\}   $. Most often the end-effector frame as seen in the world frame $ T_{ w }^{ e }  $.
It's it's own $ SE(3)  $ object in code with functions like act() which
apply the transformation to a different object.
Mathematically,
	unless explicitly stated otherwise, it's should be considered as a homogeneous 
	transformation matrix 
	\begin{equation}
	T_{ a }^{ b }= 
	\begin{bmatrix} 
			R_{ a }^{ b } \in \mathbb{R}^{3 \times 3} & p_{ a }^{ b } \in \mathbb{R}^{3} \\
			0 \in \mathbb{R}^{1 \times 3} & 1
	\end{bmatrix} 
	\end{equation}
\item $ J_{ a }^{ b } (q)  $ - Jacobian of the $ T_{ a }^{ b } (q)  $ map
	which depends on the joint configuration $ q  $.
	Most often $ J_{ w }^{ e } (q)  $ and abbreviated to just $ J  $,
	but should be obvious what it is within the context.
	\textbf{VERY IMPORTANT NOTE} by default, in Pinocchio,
	and in SMC, all Jacobians
	are \textbf{body Jacobians}, meaning they map to velocities
	in the end-effector frame $ \left\{ e \right\}   $,
	\textbf{not} in the world frame $ \left\{ w \right\}   $.
	This is because Pinocchio does kinematics and dynamics with the 
	Lie group approach in which body frame are more natural, 
	and because more computationally efficient to
	compute quantities in the body frame.
	Space Jacobians can be computed easily by passing an additional argument
	to the pinocchio::getFrameJacobian() function,
	but this is not done in the SMC framework \textit{at all}.
\item $ e  $ - generic error vector 
\item twist $ \mathcal{V}  \in \mathbb{R}^6$- combined vector of linear $ v  $
	and rotational velocity $ \omega  $, 
	$ \mathcal{V} = \begin{bmatrix} v & \omega \end{bmatrix}^{ T }  $
\item wrench $ \mathcal{F} \in \mathbb{R}^6 $- combined vector of linear force $ f  $
	and torque $ \tau  $, 
	$ \mathcal{F} = \begin{bmatrix} f & \tau \end{bmatrix}^{ T }  $
\end{itemize}
